\begin{textsample}{POS Dim 3 – persona\_grok – Score 20.00 – t932\_grok.txt}  \label{ex:f3_pos_031}
The World Health Organization ( WHO ) classifies processed \textbf{meat} as carcinogenic and red \textbf{meat} as probably carcinogenic, which raises significant concerns about current \textbf{levels} of \textbf{meat} \textbf{consumption}.
Producing the \textbf{meat} and \textbf{dairy} for just one person 's annual \textbf{needs} requires 403,000 liters of water, an \textbf{amount} \textbf{equivalent} to taking 6,190 showers.
In 2011, 8,481 \textbf{tons} of antibiotics were sold for \textbf{use} in EU \textbf{livestock}.
Animal \textbf{agriculture} involves breeding 1 \textbf{billion} pigs, 19 \textbf{billion} chickens, and 1.4 \textbf{billion} \textbf{cattle}, and it is responsible for 80 % of deforestation in the Amazon as well as emitting more \textbf{greenhouse} \textbf{gases} than the entire \textbf{transport} \textbf{sector}.
This \textbf{industry} also \textbf{uses} 80 % of global soy \textbf{production} for animal feed.
To address these issues, consider reducing your weekly \textbf{meat} intake by \textbf{half} and incorporating more local fruits and vegetables, which can offer benefits for \textbf{health}, ethics, and the environment.

% matched lemmas: agriculture, amount, billion, cattle, consumption, dairy, equivalent, gas, greenhouse, half, health, industry, level, livestock, meat, need, production, sector, ton, transport, use
\end{textsample}
